\section{Conclusion}\label{sec:conclusion}

A threshold rule for a scalar resource can be certified rather than simulated. The certificate is a robust viability kernel of a registered map over a declared class of driver floors and residual defects. In one dimension its boundary is computed by backward recursion; at infinite horizon the kernel is the safe set, a half-line above the smallest fixed point of the worst-floor map, or empty after a finite horizon. A staged rule affects a threshold only through the control active there, and a threshold is securable by some admissible rule exactly when the most protective rule's worst-floor fixed point clears it. For affine closed-loop maps the exact bounded-defect boundary coincides with the Lipschitz certificate wherever the backward recursion, rather than the inflated initial-state constraint, binds.

For Northern cod, the LRP is self-viable under constant catches up to $91.6\kt$ under the q10 floor, but under no catch under q05 or the worst floor. The Stage-I rule supplies less than the flat 60-kt cap at equal protection and is not retained; this no-dominance verdict does not depend on the unidentified carrying capacity. Exact horizons range from 8--13 years under the worst/q05 floors and 12--28 years under q10, compared with 5--7 under the registered Lipschitz certificate. For the Edwards Aquifer, a cut of at least 7.2\% secures 618 ft; Stage I and the four-stage cascade increase supply relative to their matched flat caps at comparable protection. The positive-pumping horizon is three years on the declared defect and two on the audited defect; zero pumping adds one year in each case. No admissible rule secures the 660-ft trigger under any declared floor.

The joint lesson is that the exact kernel can materially extend cod horizons because the Schaefer map's slope falls away from the LRP, while the aquifer's local zero-pumping and staged-rule departures change boundary locations but not the audited horizons. The two locked panels, archived campaign artefacts, independent core and verifier allow the results to be re-derived and cross-checked; tables and figures are regenerated from the core. Re-certification should follow the shortest horizon supported by the model class and the observed out-of-sample audit.
